\begin{answer}
\begin{itemize}

\item \textbf{Dot-product kernel} $K(x,z) = x^Tz$: Works \textit{somewhat}. This is equivalent to the standard (non-kernelized) linear perceptron operating in the original feature space. It can only learn a linear decision boundary, so it performs well only when the data is approximately linearly separable.

\item \textbf{RBF kernel} $K(x,z) = \exp\!\left(-\|x-z\|^2 / (2\sigma^2)\right)$: Works \textit{well}. The RBF kernel implicitly maps data to an infinite-dimensional feature space where complex, non-linear decision boundaries are representable. This allows the kernelized perceptron to separate classes that are not linearly separable in the original space.

\item \textbf{Non-PSD kernel} $K(x,z) = \mathbf{1}[x = z] \cdot (-1)$: \textit{Completely fails}. This function does not correspond to any valid inner product space (its kernel matrix is $-I$, which is negative definite). As a result, the kernel perceptron does not have meaningful geometric intuition and produces essentially random predictions that do not correspond to any valid decision boundary.

\end{itemize}
\end{answer}
